\chapter*{Resumen}
\addcontentsline{toc}{chapter}{Resumen}

Este libro presenta una aproximación comparativa a la sintonización de controladores
PID en sistemas de generación aislados, tomando como caso de referencia una microred
híbrida hidro-solar ubicada en Limbani, Puno. El punto de partida es la limitación de
los métodos clásicos de ajuste cuando deben actuar sobre plantas con dinámica no lineal,
efecto de golpe de ariete, baja inercia y perturbaciones vinculadas con cambios bruscos
de carga o variaciones de irradiancia solar.

A partir de ese contexto, se desarrolla una formulación del problema de sintonización como
una tarea de optimización, donde las ganancias $K_p$, $K_i$ y $K_d$ del controlador se
buscan dentro de un espacio restringido mediante criterios de desempeño como sobreimpulso,
tiempo de asentamiento e índices integrales de error, con especial atención al ITAE. Sobre
esa base se revisan los fundamentos del enfoque metaheurístico y se presentan varias
estrategias de búsqueda, destacando el algoritmo Particle Swarm Optimization (PSO) por
su equilibrio entre simplicidad, costo computacional y calidad de convergencia.

El desarrollo muestra que una sintonización optimizada permite amortiguar con mayor
eficacia las oscilaciones de frecuencia, reducir el esfuerzo del servomecanismo hidráulico
y mejorar la robustez del sistema frente a perturbaciones operativas. Además, se discuten
las posibilidades de extender el análisis a otros algoritmos de enjambre, esquemas híbridos
y entornos de validación más cercanos a la implementación industrial.

\bigskip
\noindent\textbf{Palabras clave:} metaheurísticas, sintonización PID, microredes aisladas, optimización, PSO.
